\chapter{引言}\label{chap:intro}
大语言模型（large language model，LLM）已经席卷全球。它们能够用人类自己的语言与我们交流，也能够对我们所处的世界进行推理，因此成为写作、回答问题以及完成许多其他任务时广受欢迎的工具。

它们给出的回答，已经在影响我们的决策与行动。如果让大语言模型自己直接在世界中采取行动，会发生什么？如果让大语言模型彼此交互，又会发生什么？智能体大语言模型（agentic LLM）将大语言模型的强大能力与在现实世界中行动的能力结合起来。本书将讨论智能体大语言模型如何运作、如何行动、如何交互，以及这些能力对身为用户的我们意味着什么。

\section{为什么写这本书？}
人工智能（artificial intelligence，AI）正在改变世界。能够{\em 行动}的大语言模型，则以一种新的方式改变着世界。要理解这些变化的影响有多大，人工智能领域的学生和研究者既需要了解技术本身，也需要了解正在涌现的应用。有了这样的认识，我们就能够影响这些技术的发展方向，或许还能为其引入保障措施；更具体地说，这也有助于我们进一步推动这一领域的发展。

本书是为我们在大学开设的智能体大语言模型研究生课程编写的。智能体大语言模型汇集了人工智能研究的两个领域：智能体研究与大语言模型研究。智能体正在改变大语言模型领域；反过来，大语言模型也在改变智能体领域。

虽然介绍大语言模型及其工作原理的优秀著作已经不少，例如 \citet{jm3}，但行动能力使智能体大语言模型成为一个独立的类别。我们认为，有必要开设一门课程、编写一本书，解释智能体大语言模型如何运作，以及它们如何使以往无法实现的事情成为可能。

\section{概念与定义}
大语言模型的语言能力令人印象深刻，它们说起话来与我们十分相像。如今，大语言模型已经具备智能体能力，其活动范围也远远超出了对话：它们能够{\em 推理}，能够{\em 使用工具}，例如编写和运行代码，能够浏览网页以检索信息，还能够与其他智能体或外部世界{\em 交互}。

由此，我们给出本书使用的两个定义。
\begin{quotation}
\noindent
{\em 智能体（agent）}执行动作，观测这些动作的效果，并利用从世界中获得的反馈来调整自己的动作策略（action policy）。\\
\\
{\em 智能体大语言模型}是整合了大语言模型能力的智能体，例如利用自然语言交流、编写代码或处理信息的能力。
\end{quotation}

\subsection{提供协助的智能体与彼此交互的智能体}

行动能力增强了大语言模型在社会许多领域中的重要性与适用性。智能体大语言模型有哪些用途？本书将其主要应用场景分为两类：

(1) 智能体大语言模型用于{\em 协助}人类用户，应用领域包括医疗、物流和金融等（第四部）。

(2) 智能体大语言模型用于彼此{\em 交互}，以完成任务，或在模拟人类行为的系统中发挥作用（第五部）。

近年来，这两个方向都发表了大量研究论文。本书将在各处介绍其中的关键研究与发展趋势。

\subsection{智能体控制框架}
从工程角度来看，我们该如何构建智能体大语言模型？智能体大语言模型由两个要素组成：(1) 大语言模型；(2) 一段称为{\em 控制框架（harness）}的代码。
\begin{center}智能体大语言模型 = 大语言模型 + 控制框架\end{center}
控制框架是一段包裹模型的计算机代码，用来控制调用模型的方式。英文 harness 原指马具；正如马具能够帮助我们引导骏马的力量，使之用于有益的工作，控制框架也起着类似的作用。

控制框架承担提示构造、回答处理、外部工具调用、状态管理、系统提示管理和技能管理等任务。它反复执行完成任务所需的推理、行动与观测过程，无须人类介入并手动引导。因此，智能体控制框架（agentic harness）可以看作一名用途广泛的自动提示工程师（第二部，第~\ref{chap:reasoning}~章）。尤其值得注意的是，控制框架还可以利用强化学习（reinforcement learning，RL），根据智能体动作所获得的反馈，{\em 调整}其行为策略（behavior policy）（第三部，第~\ref{chap:agents}~章）。

\subsection{多模态模型}
对于需要与用户及周围世界交互的智能体而言，理解自然语言至关重要。不过，许多任务还要求智能体处理其他模态，例如视觉信息，或其他形式的文本，例如电子表格与计算机代码。因此，智能体大语言模型的范围超出了自然语言，还包括多模态模型（multimodal model）、代码模型与世界模型（world model）。
请注意，本书根据实际语境使用{\em 大语言模型}一词：它既可能指仅处理语言的系统，也可能具有更广的含义，即{\em 基于 Transformer、能够处理自然语言、计算机代码、图像、工具命令或其他相关模态的大型模型}（第二部，第~\ref{chap:world}~章）。当这种区别至关重要时，我们会明确指出所采用的是哪一种定义。

\subsection{训练数据}
智能体大语言模型标志着大语言模型工程又一次范式转变，而这次转变与机器学习有关。长期以来，大语言模型领域谈到学习时，重点一直是模型训练期间的学习；如今，在查询模型的推理阶段进行学习，正发挥越来越大的作用。

\begin{figure}
    \centering
    % 依照源码包 figures/virtuous5.png 重绘；中心 LLM 无连接线。
% 单向：多步→推理，社会能力→交互；其余绿色节点与相应红色节点双向相连。
\resizebox{\linewidth}{!}{%
\begin{tikzpicture}[
    x=1cm,y=1cm,
    every node/.style={font=\normalsize,align=center,inner sep=5pt},
    category/.style={draw=red!85!black,line width=1.5pt,fill=white,minimum height=.58cm},
    concept/.style={draw=green!65!black,line width=1.5pt,fill=white,minimum height=.58cm},
    flow/.style={-{Stealth[length=2.6mm,width=2mm]},line width=.85pt},
    datalabel/.style={font=\small,fill=white,inner sep=2pt}
]
\node[category] (reason) at (7.1,4.1) {推理};
\node[category] (act) at (9.2,2.0) {行动};
\node[category] (interact) at (4.9,2.0) {交互};
\node[draw=black,line width=1.5pt,font=\bfseries\large,fill=white,minimum height=.78cm] (llm) at (7.1,2.8) {LLM};
\node[concept] (multi) at (4.1,5.7) {多步};
\node[concept] (reflect) at (7.1,6.2) {自我反思};
\node[concept] (retrieve) at (9.8,5.7) {检索};
\node[concept] (world) at (12.1,3.7) {世界模型、VLA};
\node[concept] (tools) at (12.1,1.8) {机器人、工具};
\node[concept] (assist) at (10.5,.25) {助手};
\node[concept] (social) at (1.6,3.6) {社会能力};
\node[concept] (role) at (1.6,1.8) {基于角色的交互};
\node[concept] (society) at (2.4,.25) {开放式社会};

% 三种主要类别形成有向循环；每条边生成新数据。
\draw[flow] (reason.east) to[out=-10,in=90] node[datalabel,pos=.5,above right] {新数据} (act.north);
\draw[flow] (act.south west) to[out=220,in=-40] node[datalabel,pos=.5,below] {新数据} (interact.south east);
\draw[flow] (interact.north) to[out=100,in=180] node[datalabel,pos=.52,above left] {新数据} (reason.west);

\draw[flow] (multi.south east) -- (reason.north west);
\draw[flow] ([xshift=-2mm]reason.north) to[out=135,in=225] ([xshift=-2mm]reflect.south);
\draw[flow] ([xshift=2mm]reflect.south) to[out=-45,in=45] ([xshift=2mm]reason.north);
\draw[flow] (reason.north east) .. controls (7.95,5.4) and (8.6,5.85) .. (retrieve.west);
\draw[flow] (retrieve.south) to[out=-90,in=0] (reason.east);

\draw[flow] (act.north east) to[out=100,in=180] (world.west);
\draw[flow] (world.south west) to[out=-90,in=0] (act.north east);
\draw[flow] (act.east) to[out=50,in=130] (tools.north west);
\draw[flow] (tools.south west) to[out=230,in=-50] (act.east);
\draw[flow] (act.south) to[out=0,in=90] (assist.north);
\draw[flow] (assist.west) to[out=180,in=-90] (act.south west);

\draw[flow] (social.south east) -- (interact.north west);
\draw[flow] (role.north east) to[out=60,in=150] (interact.west);
\draw[flow] (interact.south west) to[out=250,in=-30] (role.east);
\draw[flow] (society.north east) to[out=110,in=180] (interact.south west);
\draw[flow] (interact.south) to[out=-70,in=0] (society.east);
\end{tikzpicture}%
}

    \caption{连接三类智能体大语言模型的递归式自我改进循环：推理、行动与交互（红色，分别对应第二部、第三／四部与第五部）。影响各类别的概念以绿色表示。}
    \label{fig:virtuous}
\end{figure}

这有两个优点。第一，它为动态的智能体行为提供了灵活性与空间，第~\ref{chap:meta}~章和第~\ref{chap:agents}~章将进一步介绍这一点。第二，如今几乎整个互联网都已经被用作模型预训练的数据 \citep{sutskever2024}，因此，生成更多相关且高质量的数据就变得很有价值。通过与世界交互，智能体能够自行生成这样的新数据。图~\ref{fig:virtuous}~展示了这一过程。这些交互数据可以用于进一步训练。

视觉语言动作模型（vision-language-action model，VLA）从机器人动作序列中学习，就是大语言模型利用自身智能体动作进行训练的一个例子（第~\ref{chap:world}~章）。另一个例子见第~\ref{chap:emergent}~章：智能体社会使模型能够通过交互变得更智能，就像人类在彼此交往时能够学到新的行为与观念一样。

因此，除了支持有用的应用之外，智能体大语言模型的另一个驱动力，是生成更多训练数据的机会，由此形成递归式自我改进（recursive self-improvement）循环。图~\ref{fig:virtuous}~展示了这一通过交互生成数据的良性循环。红色表示主要类别：推理（第二部）、行动（第三／四部）和交互（第五部）。绿色表示影响某一类别的概念，这种影响既可能是直接的，也可能通过反馈循环产生。

\section{读者与目标}
本书主要面向人工智能与计算机科学专业的研究生和研究者。其他领域的学生与学者，例如认知科学、社会科学或人文学科的读者，只要有主动钻研机器学习与技术挑战的兴趣，也可能会对本书感兴趣。

本书旨在介绍并回顾智能体大语言模型领域的技术与应用进展。我们将较深入地讨论智能体大语言模型工程中的主要概念，提供科学文献的参考信息，以及源代码和数据集的链接，方便读者动手尝试构建自己的智能体。需要注意的是，这一领域还很年轻，仍在迅速变化。要跟上最新进展，深入理解基础知识至关重要。我们的目标，是兼顾基础与前沿。

我们使用本书讲授莱顿大学（Leiden University）的“智能体大语言模型”课程，该课程属于计算机科学硕士项目的人工智能方向，以及创意智能与技术（Creative Intelligence \& Technology）硕士项目。配套网站提供讲义幻灯片、作业和新闻等辅助材料。\footnote{\url{https://agenticllms.liacs.nl}}

亲爱的读者，希望你能像我们享受写作过程一样，享受阅读本书的过程。智能体大语言模型研究是一个新领域，本书同样是一部新作。新的进展每天都在发生。我们非常欢迎你通过网站提出反馈和建议，帮助我们改进本书的后续版本！

\section{预备知识}
本书假定你已经具备计算机编程、人工智能、机器学习和自然语言处理方面的本科层次知识。

我们假定你会使用 Python 和 PyTorch，也了解机器学习的基本原理，例如分类与回归的区别、训练集与测试集的区别，以及过拟合与欠拟合的成因。你应当了解决策树、集成方法和随机森林等经典机器学习方法，也应当熟悉深度学习中的一些概念，包括人工神经网络的架构（全连接多层感知机、卷积神经网络和循环神经网络）、通过反向传播进行学习、损失函数（如交叉熵）与激活函数（如 ReLU）。我们还希望你熟悉强化学习：知道什么是序贯决策问题，理解智能体与环境的区别，并了解如何使用一些策略学习算法，例如 Q-learning、DQN、PPO 和 SAC。最后，你应当对自然语言处理有所了解：什么是词元化（tokenization），什么是词嵌入（word embedding），什么是注意力（attention），以及什么是 Huggingface。

不过，我们也认为，即使你对上述部分乃至全部主题和概念了解不深，阅读本书仍然可能有所收获。我们会简要介绍必要的知识，使你能够对智能体大语言模型获得有实质内容的理解，并找到进一步学习的入口。对于其中一些主题，你或许可以在手边备一本合适的教材，\footnote{参见 \citet{gunasekar2024textbooks}。}例如：
\begin{itemize}
    \item 经典机器学习：\citet{alpaydin2021machine,flach2012machine}；
    \item 深度监督机器学习：\citet{bishop2023deep,goodfellow2016deep,chollet2021deep,stevens2020deep}；
    \item 强化学习：\citet{sutton2018reinforcement,graesser2019foundations,plaat2022deep,franccois2018introduction}；
    \item 自然语言处理与 Transformer：\citet{jm3}。
\end{itemize}

\section{课程组织方式}
我们将本书用于大学开设的智能体大语言模型研究生课程。写作时，我们希望兼顾简明与全面，以适合研究生课程的深度覆盖这一领域的完整范围。为了方便读者深入学习，我们提供了大量科学文献线索，并在带有 * 标记的章、节中介绍一些进阶主题。

不同大学的课程项目可能选取不同的内容（见表~\ref{tab:flow}）。我们开设的是综合课程，共有 14 次讲课：首先是一次导论，然后分别用 3 次讲课介绍第一部（语言建模），用 3 次讲课介绍第二部（推理）与第三部（智能体），用 3 次讲课介绍第四部（工具与助手），再用 3 次讲课介绍第五部的社会模拟（social simulation）。最后以一次综合评估课结束。我们将四个案例研究作为课程作业，并通过与学生面对面交流的方式评定成绩。课程结束时不另设理论考试，课程成绩取各次作业成绩的平均值。

当然，你可以根据自身情况，以任何合适的方式使用本书；我们重点讨论智能体大语言模型的两类主要应用场景：(1) 助手；(2) 社会模拟。

如果你的培养方案已经包含一门深入介绍大语言模型的独立课程，那么，以 (1) 智能体助手为重点的核心课程可以只选用第二、三、四部。作业方面，我们建议采用案例研究 2 和 3。

另一种安排是，以 (2) 社会模拟为重点的核心课程应当涵盖第二、三、五部，并将案例研究 2 和 4 作为作业。

无论采用助手方向还是社会模拟方向，都可以根据需要补充第一部（语言建模）；具体取决于所在大学培养方案的结构。

第六部包含进阶主题。介绍合成数据（synthetic data）生成的第~\ref{chap:synthetic-data}~章适合与第四部的助手应用场景配合使用。介绍机制可解释性（mechanistic interpretability）的第~\ref{chap:mi}~章，适合在第一、二部语言模型与推理模型的基础上深入学习。介绍先进架构的第~\ref{chap:arch}~章同样适合补充第二部。介绍认知架构（cognitive architecture）的第~\ref{chap:cognition}~章，则适合与第五部的社会模拟应用场景配合使用。

\begin{table}
    \centering
    \begin{tabular}{l|l|l|l|l}
        课程目标 & 分部 & 案例研究 & 章节 & 选修章节 \\ \hline
        综合课程 & I, II, III, IV, V & 1, 2, 3, 4 & 1-18 & 19-22\\
        助手核心课程 & II, III, IV & 2, 3 & 6-14 & 2-5, 19\\
        社会模拟核心课程 & II, III, V & 2, 4 & 6-10, 11, 15-18 & 2-5, 22\\
    \end{tabular}
    \caption{不同课程目标与本书分部、章节的对应关系}
    \label{tab:flow}
\end{table}

表~\ref{tab:flow}~概括了围绕三种课程目标使用本书材料时，可以采用的不同方案。

\section{方法、范围与全书概览}
本书强调动手实践。对于每一个新概念，我们既介绍理论背景，也举例说明如何将其用于实践。我们讨论数据集、训练算法、基准测试和训练环境，给出重要智能体算法的伪代码，并提供实际案例研究，帮助读者学习。

本书各章分别围绕一个完整主题展开，并组织为七个分部。

第一部介绍大语言模型的基础知识。第~\ref{chap:llm}~章讨论重要的自然语言处理概念，包括词元化、嵌入（embedding）、循环神经网络、注意力机制、Transformer 架构与自回归训练（autoregressive training），其中自回归训练是大语言模型训练的核心机器学习方法。我们也会进一步讨论监督学习和强化学习。接下来，第~\ref{chap:pipeline}~章较详细地介绍用于训练大语言模型的训练流程，讨论数据集获取、预训练（pretraining）、微调（finetuning）、指令微调（instruction tuning）、价值对齐（value alignment）和模型推理。我们还会介绍通常用于测试大语言模型性能的基准测试。第~\ref{chap:scaling}~章讨论扩展定律（scaling law）、小语言模型以及能源消耗。要使用智能体大语言模型，就必须理解 Transformer 模型预训练与微调的基础知识。在第一部的最后一章，我们通过一个实践案例研究，指导你从零开始训练 Transformer 模型，帮助你获得这方面的经验。

第二部介绍推理、反思（reflection）与世界模型。这三个概念都是（高级）智能体大语言模型的重要组成要素。第~\ref{chap:reasoning}~章从上下文学习（in-context learning，ICL）开始，讨论思维链（chain of thought，CoT）方法如何使大语言模型能够对简单的数学文字题进行多步推理。我们还会讨论如何借助控制框架，让大语言模型对自身的上下文学习进行自我反思（self-reflection），以增强推理能力。到这里为止，我们关注的都是使用语言的模型。要在现实世界中行动，智能体还必须理解图像数据及其他模态。第~\ref{chap:world}~章将超越语言，讨论能够理解视觉与语言的模型，以及能够从决策序列中学习并掌握机器人动作的模型。为了在世界中行动而不撞到墙上，智能体需要对物理规律有所理解。第~\ref{chap:world}~章将讨论世界模型，以及它们与基于模型的强化学习（model-based reinforcement learning）的关系。此时，我们已经讨论了智能体学习以及与环境交互的多种不同方式。第~\ref{chap:meta}~章将进一步深入讨论它们如何学习，以及这种学习为何是一种元学习（meta-learning）。

接下来的第三部中，第~\ref{chap:agents}~章从人工智能的视角回顾智能体的发展历史，并讨论智能体大语言模型的不同智能体交互方式。此外，为了在世界中行动，智能体还需要某种形式的记忆，本章将介绍不同的智能体记忆架构。第三部以一个实践案例研究结束，你可以在其中构建一个能够自我反思的智能体，挑战 BALROG 基准测试套件中的谜题与游戏。

第四部将重心从理论概念转向实际实现。我们将讨论研究社区为实现智能体大语言模型、使其能够使用工具在世界中行动而开发的框架。其中有些是专有框架，但许多框架是开放的，可以从 Github 下载。这一部讨论能够采取行动、为用户完成实际任务的智能体。第~\ref{sec:tools}~章介绍为实现智能体大语言模型而开发的工具链、框架与应用程序编程接口（application programming interface，API）。第~\ref{sec:rag}~章介绍检索增强方法，以及其他用于搜索最新信息的方法；这些信息在大语言模型训练时尚不可得。第~\ref{sec:assist}~章概览智能体大语言模型的实际应用，包括购物助手、医疗助手、交易助手以及帮助开展科学研究的助手。第四部的最后一章提供一个实践案例研究，介绍如何构建医疗诊断智能体；这也将帮助你理解它的不足，使你永远不会盲目遵从智能体给出的医疗建议。

第五部采用另一种视角：从行为与认知的角度来考察智能体大语言模型，而不以技术为着眼点。首先，我们定义并概括介绍机器行为（machine behavior）研究；随后，第~\ref{chap:behavior}~章讨论如何在受控实验环境中研究单个大语言模型与智能体大语言模型，以测试它们的社会行为。第~\ref{chap:role}~章转向多智能体系统（multi-agent system，MAS），先介绍基于角色的交互（role-based interaction）：通过构建智能体大语言模型，或向其提供提示，使其承担特定角色，并围绕目标进行交互。接着，第~\ref{chap:emergent}~章转向由多个智能体组成、以开放方式交互的系统，研究其中的{\em 涌现行为（emergent behavior）}，例如合作与规范执行。本部的案例研究将使你深入体验智能体的社会行为。

接下来的第六部包含更进阶的内容。第~\ref{chap:synthetic-data}~章讨论合成数据生成。与所有深度学习方法一样，大语言模型在很大程度上仍是一个黑箱；我们将讨论一些能够揭开其部分面纱的技术。第~\ref{chap:arch}~章讨论先进架构，包括在决策轨迹上训练的 Transformer，以及循环式与递归式大语言模型。我们也会讨论先进智能体架构：上下文强化学习、能够自我修改的智能体与元智能体。第~\ref{chap:mi}~章重点介绍机制可解释性工具，用以描绘 Transformer 架构中神经元活动的模式，从而更好地理解大语言模型的推理过程。

第~\ref{chap:cognition}~章回顾有关智能体大语言模型行为与认知的近期文献，包括心智理论（theory of mind，ToM）与意识（consciousness）等进阶问题。

最后，第七部汇集本书介绍的智能体大语言模型各项要素。第~\ref{chap:concl}~章讨论智能体大语言模型目前能够做什么、仍然不能做什么，并以对其未来前景的展望结束全书。

请注意，一些章、节带有 * 标记。这些部分包含更进阶的内容，可能更适合作为学位论文的研究材料，而不是课程内容。

\section{致谢}
本书源于我们对智能体大语言模型的着迷。这份兴趣促成了我们与学生、同事之间的许多讨论，也催生了令人振奋的研究。在某个时刻，我们觉得有必要写下这本书，并开设一门研究生课程。我们相信，智能体大语言模型是一个十分重要的主题，许多人工智能专业的学生都会希望了解它。

本书的部分材料以我们早先的工作为基础，并从中得到启发，例如 \citet{van2023theory,van2024hybrid,huisman2021survey,moerland2023model,plaat2022deep,plaat2024reasoning, plaat2025agentic,van2024llamea,chen2025does,van2025baba,xing2025cocomposer}。我们感谢这些工作的所有共同作者；没有你们，就不可能写成这本书！

第~\ref{chap:generative_agents}~章关于生成式智能体社会的案例研究，由 Dorota Moravčiková 和 Friso van der Veen 编写与开发（基于 \citep{park2023generative}）。他们调整了自己在这一主题上已有的工作，使其适合本书。我们也感谢他们慷慨投入的时间与精力，以及他们对前面各章提出的敏锐意见。

我们感谢所在研究所的学生与同事。特别感谢 Thomas B\"ack、Joost Batenburg、Joost Broekens、Xi Chen、Lennard Froma、Bas Haring、Anna Kononova、Tom van Kouwenhoven、Thomas Moerland、Rob van Nieuwpoort、Mike Preuss、Peter van der Putten、Zhaochun Ren、Derya Soydaner、Suzan Verberne、Fien van Wetten、Annie Wong 和 Tarik Začiragić；与他们的讨论使我们获益良多。

\bibliographystyle{plainnat}
\bibliography{agenticbook,refs,case,synthetic}
